\documentclass[10pt,a4paper]{amsart}
\usepackage[margin=19mm]{geometry}
\usepackage{amsmath,amssymb,mathtools}
\usepackage[scr=rsfs,cal=euler]{mathalfa}
\usepackage{xfrac}
\usepackage[unicode,hidelinks,bookmarks=false]{hyperref}
\newcommand{\ie}{i.e.\ }
\newcommand{\cf}{cf.\ }
\newcommand{\resp}{respectively\ }
\newcommand{\Subsection}{\subsection}
\providecommand{\nobreakdash}{\nobreak\hbox{-}}
\setlength{\emergencystretch}{2em}
\multlinegap0pt
\usepackage{xcolor}
\usepackage{booktabs}
\usepackage{amssymb}
\usepackage{mathtools}
\usepackage[shortlabels]{enumitem}
\usepackage{algorithm}\usepackage{algpseudocode}
\usepackage{bm}
\usepackage{tikz}
\newtheorem{theorem}{Theorem}
\usepackage{cleveref}
\crefname{theorem}{Theorem}{Theorems}
\title{Rethinking Diffusion Models with Symmetries through Canonicalization with Applications to Molecular Graph Generation}
\author{Cai Zhou, Zijie Chen, Zian Li, Jike Wang, Kaiyi Jiang, Pan Li, Rose Yu, Muhan Zhang, Stephen Bates, Tommi Jaakkola}
\date{Source: arXiv:2602.15022v1 (CC BY 4.0)}
\begin{document}
\maketitle

\noindent\emph{Source and licence.} Rethinking Diffusion Models with Symmetries through Canonicalization with Applications to Molecular Graph Generation. Version arXiv:2602.15022v1 (CC BY 4.0), distributed by arXiv under the Creative Commons Attribution 4.0 International licence, http://creativecommons.org/licenses/by/4.0/, as recorded in the arXiv licence metadata for this exact version and shown on the abstract page https://arxiv.org/abs/2602.15022v1. Full licence notice: \texttt{licenses/arxiv-2602.15022-CC-BY-4.0.txt}.\par\smallskip

\noindent\textit{Editorial scope.} Counted unit: the article's theorem theorem (source0.tex lines 985--1001). The article's complete original proof of it is delivered with it (source0.tex lines 1003--1019, one \texttt{proof} environment of 17 lines). No mathematics is written, repaired or re-derived: the statement and the proof are the article's own lines, in the article's own notation, and no retained line is substituted or re-flowed.  Retained uncounted as the source-local context the proof needs: lines 424--441; lines 430--436; lines 438--440.  Nothing was omitted from any retained span, so the package carries no omission record.  No retained line contains a citation, so the wrapper carries no bibliography and no bibliography entry.  No retained line contains a figure environment, an image-inclusion command or drawing code, so no image is delivered and the package declares no asset dependency.  Editorial formatting only, in addition: an AMS \texttt{amsart} preamble that restates the article's own declarations and macros used by the retained lines, namely the environments \texttt{theorem} on separate counters, together with the standard packages those lines need.  The retained environments print in this document as Theorem 1 in order of appearance, whereas the article prints the same items with the numbering of its own full document.  The complete native source of the article as distributed in the arXiv e-print for this exact version is bundled unchanged as \texttt{sources/194767/source0.tex}.
\begin{theorem}[Variance decomposition under group-aligned lift]
\label{thm:canon-var-decomp}
Assume $\mathcal{G}$ acts orthogonally. Under the group-aligned lifted coupling
$G\sim\lambda$ independent of $(\widetilde Z_0,\widetilde Z_1)$ and $(Z_0,Z_1)=(G\cdot\widetilde Z_0,G\cdot\widetilde Z_1)$,
let $Z_t=(1-t)Z_0+tZ_1$ and $\widetilde Z_t=(1-t)\widetilde Z_0+t\widetilde Z_1$.
With $U:=Z_1-Z_0=G\cdot\Delta$ and $\Delta:=\widetilde Z_1-\widetilde Z_0$, we have
\begin{equation}\label{eq:canon-var-decomp}
\begin{aligned}
\mathrm{Var}(U\mid Z_t)&=
\underbrace{\mathbb{E}\!\left[\mathrm{Var}(\Delta\mid \widetilde Z_t)\mid Z_t\right]}_{\text{within-slice difficulty}}\\
&+\underbrace{\mathrm{Var}\!\left(\mathbb{E}[U\mid Z_t,G]\mid Z_t\right)}_{\text{symmetry ambiguity}\ge 0}.
\end{aligned}
\end{equation}
Consequently, invariant training without canonicalization admits larger conditional variance:
\begin{equation}\label{eq:canon-var-ineq}
\mathbb{E}\big[\mathrm{Var}(U\mid Z_t)\big]\;\ge\;\mathbb{E}\big[\mathrm{Var}(\Delta\mid \widetilde Z_t)\big].
\end{equation}
\end{theorem}

\begin{equation}
\begin{aligned}
\mathrm{Var}(U\mid Z_t)&=
\underbrace{\mathbb{E}\!\left[\mathrm{Var}(\Delta\mid \widetilde Z_t)\mid Z_t\right]}_{\text{within-slice difficulty}}\\
&+\underbrace{\mathrm{Var}\!\left(\mathbb{E}[U\mid Z_t,G]\mid Z_t\right)}_{\text{symmetry ambiguity}\ge 0}.
\end{aligned}
\end{equation}

\begin{equation}
\mathbb{E}\big[\mathrm{Var}(U\mid Z_t)\big]\;\ge\;\mathbb{E}\big[\mathrm{Var}(\Delta\mid \widetilde Z_t)\big].
\end{equation}

\begin{theorem}[Variance decomposition under group-aligned lift; \Cref{thm:canon-var-decomp} in main text]
Assume $\mathcal{G}$ acts orthogonally. Under the group-aligned lifted coupling
$G\sim\lambda$ independent of $(\widetilde Z_0,\widetilde Z_1)$ and $(Z_0,Z_1)=(G\cdot\widetilde Z_0,G\cdot\widetilde Z_1)$,
let $Z_t=(1-t)Z_0+tZ_1$ and $\widetilde Z_t=(1-t)\widetilde Z_0+t\widetilde Z_1$.
With $U:=Z_1-Z_0=G\cdot\Delta$ and $\Delta:=\widetilde Z_1-\widetilde Z_0$, we have
\begin{equation}
\mathrm{Var}(U\mid Z_t)
=
\underbrace{\mathbb{E}\!\left[\mathrm{Var}(\Delta\mid \widetilde Z_t)\mid Z_t\right]}_{\text{within-slice difficulty}}
+
\underbrace{\mathrm{Var}\!\left(\mathbb{E}[U\mid Z_t,G]\mid Z_t\right)}_{\text{symmetry ambiguity}\ge 0}.
\end{equation}
Consequently,
\begin{equation}
\mathbb{E}\big[\mathrm{Var}(U\mid Z_t)\big]\;\ge\;\mathbb{E}\big[\mathrm{Var}(\Delta\mid \widetilde Z_t)\big].
\end{equation}
\end{theorem}

\begin{proof}
Apply the standard law of total variance with respect to the latent symmetry variable $G$:
\begin{equation}
\label{eq:ltv}
\begin{aligned}
\mathrm{Var}(U\mid Z_t)=
\mathbb{E}\!\left[\mathrm{Var}(U\mid Z_t,G)\mid Z_t\right]
+\mathrm{Var}\!\left(\mathbb{E}[U\mid Z_t,G]\mid Z_t\right).
\end{aligned}
\end{equation}
The second term is nonnegative. For the first term, conditioning on $(Z_t,G)$ determines $\widetilde Z_t = G^{-1}\cdot Z_t$. Under the lifted coupling, $(\Delta,\widetilde Z_t)$ is independent of $G$, and orthogonality implies $\mathrm{Var}(G\cdot X\mid \cdot)=\mathrm{Var}(X\mid \cdot)$. Therefore,
\begin{equation}
\mathrm{Var}(U\mid Z_t,G)=\mathrm{Var}(G\cdot\Delta\mid \widetilde Z_t,G)
=\mathrm{Var}(\Delta\mid \widetilde Z_t).
\end{equation}
Substituting into~\eqref{eq:ltv} yields~\eqref{eq:canon-var-decomp}. Taking expectations over $Z_t$ and using the tower property gives~\eqref{eq:canon-var-ineq}.
\end{proof}

\end{document}
